\documentclass[10pt, letterpaper]{article}

\usepackage{../../../../template/template}

\setDocTitle{Verbale interno 24/03/2026}
\setDocVersion{-}
\setDocData{24/03/2026}
\setDocIndex{vi\_2026\_03\_24}
\setDocState{1}
\setDocRecipient{1}

\begin{document}

\makeTitlePage

% \newpage
% \addChangelog{0.1.0}{11/03/2026}{F. Pasqual}{C. Libralato}{\noOne}{\firstDraft}
% \makeChangelog

\newpage

\tableofcontents

\newpage

% -------------------------------------------------
% Informazioni riunione
% -------------------------------------------------
\makeMeetingInfoTable{24/03/2026}{13:45}{14:30}{via \textit{Discord$_G$}}

\addParticipant{Valeria}{Baleanu}{Programmatore}{P}
\addParticipant{Leonardo}{Pellizzon}{Programmatore}{P}
\addParticipant{Luca}{Granziero}{Amministratore}{P}
\addParticipant{Francesco}{Pasqual}{Programmatore}{P}
\addParticipant{Giuseppe}{De Fina}{Verificatore}{P}
\addParticipant{Christian}{Libralato}{Responsabile}{P}
\addParticipant{Filippo}{Venzo}{Programmatore}{P}
\makeMeetingParticipantsTable


\addPrecedentTodo{vi\_2026\_03\_16.a1}{\#5}{L. Pellizzon}{22/03/2026}

\addPrecedentTodo{vi\_2026\_03\_16.a1}{\#5}{L. Pellizzon}{22/03/2026}
\addPrecedentTodo{vi\_2026\_03\_16.a2}{\#6}{L. Granziero, F. Venzo}{22/03/2026}

\addPrecedentTodo{vi\_2026\_03\_16.a8}{\#208}{F. Pasqual}{20/03/2026}
\addPrecedentTodo{vi\_2026\_03\_16.a9}{\#209}{F. Pasqual}{20/03/2026}


\makePrecedentTodoTable
% -------------------------------------------------
% Ordine del giorno
% -------------------------------------------------
% -------------------------------------------------
% Ordine del giorno
% -------------------------------------------------
\section{Ordine del giorno}
\begin{itemize}
    \item Analisi struttura allarmi;
    \item strategie di aggiornamento cache dispositivi;
    \item gestione autenticazione JWT e associazione email utente;
    \item stato avanzamento sviluppo front-end e backend;
    \item organizzazione e divisione specifica tecnica;
\end{itemize}

% -------------------------------------------------
% Approfondimento
% -------------------------------------------------
\section{Approfondimento}

\subsection*{Analisi struttura allarmi}
È stata discussa la struttura dati degli allarmi con attenzione a campi quali ID univoco, activation time, resolution time, stato e associazione ad utente.
Si è deciso che lo stato non deve essere memorizzato ma calcolato in base ai valori activation time, resolution time e l'assegnazione ad un operatore al fine di evitare 
ridondanze. Inoltre si è convenuto sul fatto che il nome dell'allarme deve provenire dalla regola scatenante piuttosto che dall'evento stesso.

\subsection*{Strategie di aggiornamento cache dispositivi}
Si è affrontata la gestione della cache per i dispositivi degli impianti, che è conservata in file un JSON per ogni impianto e va aggiornata quando risulta obsoleta. 
Sono stati discussi diversi approcci per la gestione dell'età della cache e del suo aggiornamento, che consistono nell'affidamento a uno tra: \textit{cron job$_G$}, device service, adapter o servizio dedicato. 
La scelta non è ancora stata definita.

\subsection*{Gestione autenticazione JWT e associazione email utente}
È stato approfondito il flusso di autenticazione con JWT: il login causa \textit{redirect} alla pagina gialla di Vimar, che restituisce un code in \textit{callback} al \textit{backend}.
Attualmente mancano informazioni legate all'email dell'utente autenticato e allo stato di collegamento. Si è convenuto di aggiungere un endpoint "user info" che fornisca email e flag di collegamento dopo il login,
dal momento che la callback non può rispondere direttamente con queste informazioni al front-end.

\subsection*{Stato avanzamento sviluppo front-end e backend}
Il team ha illustrato lo stato attuale dello sviluppo. Sono stati discussi i test unitari con mock del backend e la necessità di mostrare almeno una demo funzionante
per il prossimo incontro con la Proponente. Sono stati valutati anche gli strumenti \textit{swagger$_G$} per documentazione API e possibili miglioramenti alla struttura degli endpoint.

\subsection*{Organizzazione e divisione specifica tecnica}
È stato proposto di dividere chiaramente la specifica in sezioni per backend e frontend, si è inoltre valutato di creare una sezione comune di \textit{DTO$_G$} condivisi fra client e servizio per evitare ridondanze. 

% -------------------------------------------------
% Decisioni
% -------------------------------------------------
\addDecision{Rimuovere il campo status memorizzato negli allarmi e calcolarlo dinamicamente.}
\addDecision{Implementare un endpoint user info che fornisca email e stato collegamento utente dopo login.}
\addDecision{Mostrare una demo nella prossima riunione con la Proponente.}
\addDecision{Dividere la specifica tecnica in sezioni distinte per backend e frontend con una sezione comune per DTO condivisi.}


\makeDecisionTable

% -------------------------------------------------
% Azioni
% -------------------------------------------------
\addTodo{\noOne}{Definire la modalità definitiva di aggiornamento cache dispositivi.}{\teamName}{31/03/2026}
\addTodo{\#24}{Creare endpoint "user info" per fornire email e flag collegamento utente dopo login con JWT.}{F. Venzo}{31/03/2026}
\addTodo{\noOne}{Integrare frontend e backend.}{\teamName}{26/03/2026}
\addTodo{\#224}{Redigere il verbale della riunione.}{C. Libralato}{31/03/2026}
\addTodo{\#225}{Aggiornare glossario.}{C. Libralato}{31/03/2026}

\makeTodoTable


\end{document}